\chapter{Future Research Directions}

%szb: nem szoktak szerintem ennek totál külön chaptert létrehozni
% ez mehet a conclusion-be, vagy a limitations-be
% vaaagy akár mindkettőbe: limitációkra reagálva részletesen kifejtve egy lehetséges future irányt + conclusion végén röviden vázlatosan összefoglalva, hogy mik a lehetséges irányok


Several directions can be explored in future work to further improve the quality
and reliability of automatic code translation systems.

First, future studies should evaluate additional programming language pairs.
This thesis focused only on Java-to-C++ translation, but other language pairs
may present different challenges due to differences in syntax, typing systems,
memory management, or programming paradigms. Recent work has shown that
translation difficulty can vary significantly depending on the selected source
and target languages, and that some directions are inherently more difficult
than others \cite{aljagthami2025prompt}.

Another promising direction is the use of larger models. While 3B parameter
models were selected in this thesis due to hardware limitations, larger code
models may achieve better functional correctness and stronger generalization.
Recent studies suggest that larger language models consistently outperform
smaller ones on code translation tasks, particularly for more complex programs
and language pairs \cite{yang2024llm, aljagthami2025prompt}.

Future work could also compare additional evaluation metrics beyond BLEU and
execution-based evaluation. Metrics such as CodeBLEU, CodeScore, chrF, or
METEOR may provide further insight into syntactic quality, semantic similarity,
and code structure. Prior work has shown that no single metric is sufficient
for capturing all aspects of translation quality, motivating the use of
multiple complementary evaluation methods \cite{evtikhiev2022bleu,
codescore2023}.

Another important research direction is prompt engineering. Recent studies have
demonstrated that prompt wording, prompt language, and prompt detail can
significantly influence translation quality. For example, detailed prompts were
shown to outperform concise prompts, while prompt language also affected the
quality of the generated translations \cite{aljagthami2025prompt}. Other work
has further shown that prompt engineering can substantially improve the
performance of large language models in code-related tasks \cite{chen2025prompt}.

Future studies may also benefit from larger and more diverse training datasets,
as well as more advanced fine-tuning strategies. For example, multi-stage
fine-tuning, retrieval-augmented generation, and dataset generation methods
that include execution feedback or intermediate reasoning steps may further
improve translation quality \cite{chen2025dialogue, zhu2022xlcost}.

Finally, future research should investigate multi-file and repository-level code
translation. Most existing studies, including this thesis, focus on isolated
functions or short code snippets. However, real-world software systems often
contain dependencies across multiple files, libraries, and modules. Recent work
suggests that repository-level translation remains a major challenge, but also
offers significant potential for improving practical software migration
scenarios \cite{zhang2025repository, bruneel2025pipeline}.